\chapter{Examples of LLM Prompts and Responses}
\label{annex:llm-prompts}

This annex illustrates the three language-model calls using one successful mission from the final minimal-contract evaluation:

\begin{lstlisting}[style=annexcode,caption={User mission supplied to all three calls.}]
Set the joint-state publication frequency to 30 hertz.
\end{lstlisting}

The responses below are the actual raw JSON objects returned by \texttt{qwen2.5-coder:7b} for case \texttt{seed-108-p01}. The system-message examples preserve the operative instructions and representative runtime context. Large generated catalogue, candidate-record, and JSON-schema sections are abbreviated because they are machine-generated inputs rather than hand-written instructions; their omission is marked explicitly.

\section{Capability Interpretation Prompt}

The first call determines the complete capability state. Parameter-only language must not change capabilities, so unmentioned capabilities retain their registry-defined defaults.

\begin{lstlisting}[style=annexcode,caption={Representative capability-interpretation system message.},label={lst:capability-system-prompt}]
You translate an operator mission into one MissionInterpretation JSON object.
Return JSON only. Never emit prose, ROS YAML, launch code, or fields absent from the schema.

Capability output is complete and minimal. Always return exactly:
- mapping: "cartographer" or null
- navigation: "nav2" or null
- exploration: "explore_lite" or null
- odometry: "kinematic_icp", "kiss_icp", or null

An implementation string means active. null means explicitly disabled.

Defaulting rules:
- Begin with the registry defaults.
- Preserve a default when its capability is not mentioned.
- Use null only when the operator explicitly disables a capability.
- Replace a default only when another supported implementation is named.
- Exploration requires mapping.

This call evaluates capability changes only. Parameters, topics, rates,
thresholds, and numeric settings are not capabilities. A parameter-only
request is a valid capability interpretation and preserves all defaults.

Use status "unsupported" only for an unsupported capability or
implementation. Use "needs_clarification" only when ambiguity prevents
a safe capability choice.

AntRobot capability catalogue:
  [generated registry projection omitted for readability]

Required JSON schema:
  [MissionInterpretation schema omitted for readability]
\end{lstlisting}

\subsection*{LLM response}

\begin{lstlisting}[style=annexcode,language={},caption={Capability interpretation returned by the model.},label={lst:capability-response}]
{
  "status": "valid",
  "capabilities": {
    "mapping": "cartographer",
    "navigation": "nav2",
    "exploration": "explore_lite",
    "odometry": "kinematic_icp"
  }
}
\end{lstlisting}

\section{Parameter Selection Prompt}

After deterministic capability realization, the second call receives only parameters belonging to the approved component scope. It selects identifiers but does not calculate values. Retrieval ranks, scores, and origins are not included in the model-visible records.

\begin{lstlisting}[style=annexcode,caption={Representative parameter-selection system message.},label={lst:selection-system-prompt}]
You select which ROS configuration values are relevant to the requested change.
Return exactly one ParameterSelectionInterpretation JSON object and no prose.

This stage answers only what should change. Never calculate new values.
Use only parameter_id values in the bounded selection context. Select every
occurrence required for a cross-component property. Do not select a related
parameter merely because it is nearby in the graph.

Use status:
- valid: at least one relevant parameter is identified
- no_change: no parameter change was requested
- unsupported: no supplied candidate represents the requested concept
- needs_clarification: a safe selection requires more information

Capability selection does not imply parameter tuning. "Build a map" does
not request changing launch.map. Never invent changes needed to make a
capability work. Never select parameters belonging to inactive components.

Stable system context:
{
  "active_components": [
    "cartographer_node", "explore", "joint_state_estimator",
    "kinematic_icp", "nav2", "rdrive_node", "rplidar_node"
  ],
  "inactive_components": ["kiss_icp", "laserscan_to_pointcloud"]
}

Bounded selection context:
{
  "included_parameter_ids": [
    "nodes.explore.planner_frequency",
    "nodes.explore.publish_rate",
    "nodes.joint_state_estimator.publish_frequency",
    "nodes.joint_state_estimator.wheel_radius",
    "nodes.rdrive_node.odom_frequency"
  ],
  "records": [
    {
      "parameter_id": "nodes.joint_state_estimator.publish_frequency",
      "current_value": 20.0,
      "value_type": "number",
      "semantic_metadata": {
        "description": "Controls the joint-state publication frequency.",
        "unit": "Hz"
      }
    }
  ]
  [remaining bounded candidate records omitted for readability]
}
\end{lstlisting}

\subsection*{LLM response}

\begin{lstlisting}[style=annexcode,language={},caption={Parameter selection returned by the model.},label={lst:selection-response}]
{
  "status": "valid",
  "selected_parameters": [
    {
      "parameter_id": "nodes.joint_state_estimator.publish_frequency",
      "relevance": "The parameter directly controls how often joint state information is published, which affects the rate at which joint states are updated.",
      "confidence": 1.0
    }
  ],
  "reason": "The request specifies a change to the joint-state publication frequency."
}
\end{lstlisting}

\section{Parameter-Value Reasoning Prompt}

The final call sees only the validated selection. It must repeat the frozen current value, derive the requested new value, and cannot introduce another identifier.

\begin{lstlisting}[style=annexcode,caption={Representative parameter-value system message.},label={lst:value-system-prompt}]
You calculate proposed values for parameters selected in the previous stage.
Return exactly one ParameterValueInterpretation JSON object and no prose.

This stage answers only how the selected parameters should change. Never
introduce a parameter absent from selected_records. Copy old_value exactly
from current_value. Perform unit conversion when sufficient information is
present. Keep coupled occurrences consistent. Give a short, externally
verifiable reason and quote only the shortest relevant user phrase as
request_evidence. Use needs_clarification instead of guessing.

Selected parameter context:
{
  "selection": {
    "status": "valid",
    "selected_parameters": [
      {
        "parameter_id": "nodes.joint_state_estimator.publish_frequency",
        "relevance": "Controls how often joint states are published.",
        "confidence": 1.0
      }
    ]
  },
  "selected_records": [
    {
      "parameter_id": "nodes.joint_state_estimator.publish_frequency",
      "current_value": 20.0,
      "value_type": "number",
      "unit": "Hz"
    }
  ],
  "relevant_relationships": []
}

Required JSON schema:
  [ParameterValueInterpretation schema omitted for readability]
\end{lstlisting}

\subsection*{LLM response}

\begin{lstlisting}[style=annexcode,language={},caption={Parameter-value interpretation returned by the model.},label={lst:value-response}]
{
  "status": "valid",
  "changes": [
    {
      "parameter_id": "nodes.joint_state_estimator.publish_frequency",
      "old_value": 20.0,
      "new_value": 30.0,
      "reason": "The request specifies a change to the joint-state publication frequency.",
      "request_evidence": "set the joint-state publication frequency to 30 hertz"
    }
  ]
}
\end{lstlisting}
